% !TEX root = task-3-23-2.tex
% Editable task bank. Append new BankTask blocks at the end of the document.
\documentclass[12pt,a4paper]{article}
\usepackage{physics-textbook}
\usepackage{graphicx,needspace}
\usetikzlibrary{decorations.pathmorphing}
\newif\ifShowSolutions
\ifdefined\TutorEdition
  \ShowSolutionstrue
\else
  \ShowSolutionsfalse % Change to \ShowSolutionstrue to display solutions.
\fi
\newcounter{banktask}
\newsavebox{\BankTaskBox}
\newcommand{\BankTaskContent}[3]{%
  \noindent\textbf{\thebanktask.}\quad #2\par
  \ifShowSolutions
    \Needspace{50mm}% Keep the solution heading with its opening text or diagram.
    #3\par
    \ifdefined\TutorEdition
      \ifcsname TutorSource#1\endcsname
        \nopagebreak\smallskip{\small\csname TutorSource#1\endcsname\par}%
      \fi
    \fi
  \fi
}
\newcommand{\BankTask}[3]{%
  \refstepcounter{banktask}\par\medskip
  \sbox{\BankTaskBox}{\begin{minipage}{\linewidth}%
    \BankTaskContent{#1}{#2}{#3}\end{minipage}}%
  \ifdim\dimexpr\ht\BankTaskBox+\dp\BankTaskBox\relax>\dimexpr\textheight-20mm\relax
    \Needspace{35mm}\BankTaskContent{#1}{#2}{#3}%
  \else
    \Needspace{\dimexpr\ht\BankTaskBox+\dp\BankTaskBox\relax}%
    \noindent\usebox{\BankTaskBox}\par
  \fi
  \medskip
}
\begin{document}
\BookSubsection{3.23.2}{Резонанс в колебательном контуре}


% Source 4508: https://phys-ege.sdamgia.ru/problem?id=4508
\BankTask{4508}{
К колебательному контуру подсоединили источник тока, на клеммах которого напряжение гармонически меняется с частотой $\nu.$ Электроемкость С конденсатора колебательного контура можно плавно менять от минимального значения $C_{\min}$ до максимального $C_{\max},$ а индуктивность его катушки постоянна. Ученик постепенно увеличивал емкость конденсатора от минимального значения до максимального и обнаружил, что амплитуда силы тока в контуре все время возрастала. Опираясь на свои знания по электродинамике, объясните наблюдения ученика.
}{
\smallskip\noindent\textbf{Решение}\par
В колебательном контуре источником тока возбуждаются вынужденные колебания. Частота этих колебаний равна частоте источника $\nu.$ Амплитуда же зависит от соотношения между внешней частотой и частотой собственных электромагнитных колебаний $\nu_0 = \frac{1}{2 \pi \sqrt{LC}} .$ По мере увеличения емкости конденсатора уменьшается частота собственных колебаний в колебательном контуре. Амплитуда колебаний достигает максимума при резонансе, когда $\nu=\nu_0.$ После этого амплитуда начинает убывать. В данном случае, ученик меняет не внешнюю частоту, а частоту собственных электромагнитных колебаний. При плавном увеличении емкости контура от минимального значения $C_{\min}$ до максимального $C_{\max},$ частота уменьшается от $\nu_0,\max = \frac{1}{2 \pi \sqrt{LC_{\min}}}$ до $\nu_0,\min = \frac{1}{2 \pi \sqrt{LC_{\max}}} .$ Из того факта, что амплитуда все время увеличивалась, можем сделать вывод, что частота $\nu_0$ все время приближалась к частоте источника тока, при этом $\nu < \nu_0,\min.$
}

\end{document}
